% ==================================================
% Difference Rule
% Discrete Calculus Notes Stub
% ==================================================

\subsection*{Difference Rule}

The falling factorial basis interacts particularly well with the
difference operator.

\begin{theorem}[Difference Rule for Falling Factorials]
For every integer \(k \ge 1\),

\[
\Delta x^{\underline{k}} = k\,x^{\underline{k-1}}.
\]
\end{theorem}

\begin{remark}
This rule is the discrete analogue of the familiar power rule from
calculus,

\[
\frac{d}{dx} x^k = kx^{k-1}.
\]

Because of this identity, falling factorials behave under the difference
operator exactly the way ordinary powers behave under differentiation.
For this reason they form the natural polynomial basis for discrete
calculus.
\end{remark}


\subsection*{Examples}

\[
\begin{array}{c|c}
f(x) & \Delta f(x) \\
\hline
x^{\underline{1}} = x & 1 \\
x^{\underline{2}} = x(x-1) & 2x \\
x^{\underline{3}} = x(x-1)(x-2) & 3x(x-1) \\
x^{\underline{4}} = x(x-1)(x-2)(x-3) & 4x(x-1)(x-2)
\end{array}
\]

\begin{remark}
In ordinary calculus the monomials \(1,x,x^2,x^3,\dots\) form the natural
basis because derivatives act simply on them. In discrete calculus the
falling factorials \(1,x,x^{\underline{2}},x^{\underline{3}},\dots\)
play the same role because the difference operator acts on them through
the simple rule

\[
\Delta x^{\underline{k}} = k\,x^{\underline{k-1}}.
\]
\end{remark}

\begin{center}
\begin{tikzpicture}
\foreach \n in {0,...,5}{
    \foreach \k in {0,...,\n}{
        \node at (\k-\n/2,-\n) {$\binom{\n}{\k}$};
    }
}
\end{tikzpicture}
\end{center}

\begin{remark}
The coefficients appearing in higher finite differences follow the same
pattern as Pascal's triangle.  In fact,

\[
\Delta^n f(x)
=
\sum_{k=0}^{n}
(-1)^{\,n-k}
\binom{n}{k} f(x+k).
\]

Thus the binomial coefficients governing Pascal's triangle also determine
the structure of higher difference operators.
\end{remark}























